\section{Resources and the limits of a single backend}
\label{sec:resources}

\subsection{Distinct correlation models}
For a fixed finite Bell arena, the unrestricted local set is
\begin{equation}
 p(a\mid x)=\sum_\lambda q(\lambda)\prod_i p_i(a_i\mid x_i,\lambda).
 \label{eq:local}
\end{equation}
Its extreme points are deterministic response functions $a_i=f_i(x_i)$.
Private randomization and shared randomization do not improve a linear optimum
beyond the best deterministic response tuple.

Finite-dimensional tensor-product quantum behaviors have the form
\begin{equation}
 p(a\mid x)=\Tr\left[\rho\bigotimes_i M^{(i)}_{a_i\mid x_i}\right],\qquad
 M^{(i)}_{a_i\mid x_i}\succeq0,\quad
 \sum_{a_i}M^{(i)}_{a_i\mid x_i}=\id.
 \label{eq:quantum}
\end{equation}
The dimension-bounded set $\Q_{\boldsymbol d}$, the union $\Q_{\rm fin}$ over
all finite dimensions, and its closure $\Q_{\rm a}$ are different types.
For a dimension bound, one must also say whether free shared classical
randomness is allowed: convexifying $\Q_{\boldsymbol d}$ can change membership,
although it does not change a linear optimum. Projective-measurement dilation
is harmless when local dimension is unrestricted, but generally changes a
fixed-dimensional resource contract.

In the commuting-operator set $\Q_{\rm c}$, all measurement operators act on a
single Hilbert space, and operators assigned to distinct parties commute.
Non-signalling behaviors satisfy, for every subset $S$ of parties,
\begin{equation}
 \sum_{a_{\bar S}}p(a_S,a_{\bar S}\mid x_S,x_{\bar S})
 \quad\text{is independent of }x_{\bar S}.
 \label{eq:ns}
\end{equation}
The useful inclusion chain is
\begin{equation}
 \Lset\subseteq\Q_{\rm fin}\subseteq\Q_{\rm a}
 \subseteq\Q_{\rm c}\subseteq\NS.
 \label{eq:chain}
\end{equation}
Equalities and strictness depend on the arena. In general $\Q_{\rm fin}$ is not
closed \cite{slofstra2019}, and $\Q_{\rm a}\ne\Q_{\rm c}$ for some finite
arenas \cite{ji2020}. The closure has the same supremum for every continuous
linear score as $\Q_{\rm fin}$, but membership and attainment can differ.
These are operational distinctions, not interchangeable solver labels.

\subsection{Three extension axes}
A useful taxonomy is by additional semantics, not by a fixed number of disjoint
classes. The following axes can occur together.

\paragraph{Independent sources and causal structure.}
A network resource specification contains a source hypergraph. In a classical
network, independent latent variables $\lambda_e$ obey
\begin{equation}
 p(a\mid x)=\int\prod_e dq_e(\lambda_e)
 \prod_i p_i(a_i\mid x_i,(\lambda_e)_{e\ni i}).
 \label{eq:network}
\end{equation}
A quantum network has a product $\bigotimes_e\rho_e$ of source states, with
local measurements on all systems arriving at each node. A global hidden
variable shared across sources destroys the independence assumption. The
resulting sets are generally nonconvex; their analysis requires causal
compatibility methods such as inflation \cite{fritz2012,wolfe2021}.
Forgetting independence is a relaxation, never a lossless normalization.

\paragraph{Trusted quantum interfaces.}
Semiquantum games use known quantum states as questions. Substituting their
classical preparation labels as freely readable inputs grants the players
information they did not possess when the states were nonorthogonal. Thus
finite labels alone do not provide a faithful Bell-game lowering. Buscemi's
semiquantum tests witness entanglement beyond ordinary Bell tests
\cite{buscemi2012}. Quantum answers, steering assemblages, and trusted referee
systems similarly require operator-valued interfaces and explicit trust types.

\paragraph{Multiple rounds and memory.}
Sequential interactions need an event order, local histories, admissible
communication, and quantum instruments or channels. A transcript table alone
does not enforce that an early answer is independent of a future question.
For finite-dimensional quantum interactions, strategy/comb representations
provide positive operators with causal trace constraints
\cite{gutoski2007}. Locality between different players remains an additional
restriction; it is not supplied by positivity of a single global comb.

The proposed umbrella IR is a typed causal process: classical or quantum wires,
trusted or untrusted nodes, independent source factors, and a referee score.
It can describe many finite protocols uniformly, but each backend must accept
only a verified fragment. Continuous interfaces, infinite rounds,
superselection rules, energy constraints, computational restrictions, and
generalized probabilistic theories add further semantics. No claim of a
complete finite classification is made.

\begin{proposition}[No-input linear-score obstruction]
\label{prop:noinput}
Suppose every party has one input and every deterministic joint output is an
admissible strategy in a resource model $\mathcal R$. Then
\[
 \val_{\mathcal R}(V)=\max_{a\in A}V(a).
\]
In particular, ordinary local and quantum models cannot be separated by such
a game. The same value obstruction applies to independent-source networks
when arbitrary constant outputs are permitted.
\end{proposition}
\begin{proof}
Every score is a convex combination of the numbers $V(a)$, hence at most
their maximum. The players attain the maximum by each outputting their
component of one maximizing tuple, without using any source.
\end{proof}
No-input network nonlocality is nevertheless meaningful: it concerns whether
a \emph{distribution} belongs to a nonconvex compatible set
\cite{fritz2012}. Its detection generally needs a nonlinear distribution
witness, constraints fixing other statistics, or several measured quantities.
Such tests must be represented as compatibility tasks rather than misleadingly
reported as an unrestricted scalar winning-probability advantage.
